\section{11. Conclusion}

\subsection{11.1 Summary}

We presented \textbf{Galerkin Attention}, which formulates attention as a 
Petrov-Galerkin projection onto a DCT basis, achieving $\mathcal{O}(Nrd)$ 
complexity for encoder models. Decoders fall back to standard softmax.

\subsection{11.2 Results Recap}

\begin{table}[H]
\centering
\begin{tabular}{llll}
\toprule
\textbf{Model} & \textbf{Acc Δ} & \textbf{Params Δ} & \textbf{Path} \\
\midrule
DistilBERT & -0.80\% & +589824M & Galerkin \\
RoBERTa & -1.21\% & +1179648M & Galerkin \\
GPT-2 & 0.00\% & +1188876M & Fallback \\
\bottomrule
\end{tabular}
\end{table}

\subsection{11.3 Key Takeaways}

1. \textbf{Encoder linear complexity} is achievable via basis projection.
2. \textbf{DCT is the natural basis} for attention smoothness.
3. \textbf{Decoder fallback} is necessary for causality.
4. \textbf{Not a compression method.}

\subsection{11.4 Future Work}

\subsubsection{Short-Term}
\begin{itemize}
\item Learned basis
\end{itemize}
\begin{itemize}
\item Causal-compatible basis
\end{itemize}
\begin{itemize}
\item Cross-attention evaluation
\end{itemize}

\subsubsection{Long-Term}
\begin{itemize}
\item Higher-order DCT variants
\end{itemize}
\begin{itemize}
\item Adaptive rank per layer
\end{itemize}
\begin{itemize}
\item Integration with FlashAttention
\end{itemize}

\subsection{11.5 Closing Remarks}

Galerkin Attention is a mathematically principled encoder efficiency method. 
We release code under MIT.

\subsection{11.6 Acknowledgments}

{{ACKNOWLEDGMENTS}}

\subsection{11.7 Author Contributions}

\begin{itemize}
\item \textbf{{{AUTHOR_1}}} designed the method, implemented the code, ran experiments.
\end{itemize}
\begin{itemize}
\item \textbf{{{AUTHOR_2}}} provided feedback.
\end{itemize}

\subsection{11.8 Conflict of Interest}

None.